\centerline{\Huge \bf  Экзаменационная работа по алгебре 8 класс}
\centerline{\Huge \bf курс углубленного изучения математики}
\centerline{\Large \bf 2 вариант}

\begin{flushright}
    \scriptsize{Логика может привести Вас от пункта А к пункту Б,\\ а воображение — куда угодно.\\ Альберт Эйнштейн.}
\end{flushright}

\problem{1}
Упростите выражение
\[ \left( \dfrac{1}{a^2 - x^2} + \dfrac{x-a}{x+a}\right)\cdot\dfrac{x-a}{x^2 - a^2 + x+a} + \dfrac{2a+1}{(x+a)^2}\]
\problem{2}
Между какими целыми числами находится каждый корень уравнения

\noindent
$\frac{2}{2x+5} - \frac{2x+1}{4x^2-25} + \frac{1}{6} = 0$

\problem{3}
Решить уравнение\\
$(x^2-7x+13)^2 - (x-3)(x-4) = 1$

\problem{4}
\textbf{Часть 1.} Постройте график функции $f(x) = \sqrt{x^2-4x+4}-3$ и укажите: а) нули функции; б) область определения; в) область значений функции. 

\noindent
\textbf{Часть 2.} Постройте график функции $g(x) = \dfrac{x-2}{x-2} + \dfrac{2}{x+1}$ и укажите: а) нули функции; б) область определения; в) область значений функции.

\noindent
\textbf{Часть 3.} Укажите число точек пересечения графиков $f(x)$ и $g(x).$

\problem{5}
Найдите последнюю цифру числа: 
$17^9 - 23 \cdot 2^7 - 3^{15}$

\problem{6}
Докажите равенство:
\[
\left( \frac{11}{5 - \sqrt{3}} \right)^2 - \left( \frac{5 - 2\sqrt{5}}{2 - \sqrt{5}} \right)^2 = \sqrt{ \frac{91}{4} + 10\sqrt{3} }.
\]

\problem{7}
Выяснить, какое наименьшее значение может принимать сумма квадратов корней уравнения $x^2 - (2a-3x) + a^2 - 3a = 0$

\problem{8}
\text{Решить совокупность неравенств:}
$
\left[ 
  \begin{gathered} 
    \left\{ 
      \begin{gathered} 
        x+1>\dfrac{x-1}{4}
        \\  
        \dfrac{x-1}{3}\geq \dfrac{x+1}{5}-\dfrac{1}{15}
        \\
      \end{gathered} 
    \right.  
    \\ 
    |x-2| < 1,5
  \end{gathered}
  \right.
$

\problem{9}
Из пункта А в пункт В, расстояние между которыми 80 км, отправился велосипедист. Одновременно навстречу ему из пункта В отправился мотоциклист. Велосипедист прибыл в пункт В через $3$ ч после их встречи, а мотоциклист прибыл в пункт А через $1$ ч $20$ мин после встречи. Определите, на каком расстоянии от пункта А произошла встреча?


\newpage

\hspace{3mm}

\problem{10}
Найдите натуральные числа $k$ и $l$, если известно, что из трех следующих утверждений два истинны, а одно --- ложно:

1) $5k+8l=120;$ \qquad 2) $8k+5l = 120$ \qquad 3) $7k+10l=195.$ 